% Môn: Toán
% Lớp: 12
% Chương: Chương I. Ứng dụng đạo hàm để khảo sát và vẽ đồ thị hàm số
% Bài: Bài 1. Tính đơn điệu và cực trị của hàm số
% Số câu: 16
% App đọc file này trực tiếp — sửa rồi Commit trên GitHub, bấm Đồng bộ GitHub .tex

% Bài 1. Tính đơn điệu và cực trị của hàm số

% ===== Câu 1 =====
% ID: T12C1B1-01-DS
% Mức: TH
\begin{ex}
% ID: T12C1B1-01-DS
% Mức: TH
$\immini$
 {
 Cho hàm số $y=f(x)$ có đồ thì như hình vẽ bên. Các mệnh đề sau đúng hay sai?
\choiceTF
{\True Hàm số đã cho nghịch biến trên khoảng $(-2,0)$}
{Hàm số đã cho đồng biến trên khoảng $(-3;1)$}
{\True Hàm số đã cho đồng biến trên khoảng $(2;+\infty)$}
{Hàm số $f=|f(x)|$ nghịch biến trên khoảng $(-2,0)$}
\loigiai{
\begin{itemchoice}
\item (Đúng) Hàm số đã cho nghịch biến trên khoảng $(-2,0)$. (Vì): Trên khoảng $(-2,0)$, đồ thị hàm số đi xuống từ trái sang phải, nên hàm số đã cho nghịch biến trên khoảng này
\item (Sai) Hàm số đã cho đồng biến trên khoảng $(-3;1)$. (Vì): Trên khoảng $(-3;1)$, đồ thị hàm số có cả đoạn đi lên (đồng biến) và đoạn đi xuống (nghịch biến), nên hàm số không đồng biến trên toàn bộ khoảng này
\item (Đúng) Hàm số đã cho đồng biến trên khoảng $(2;+\infty)$. (Vì): Trên khoảng $(2;+\infty)$, đồ thị hàm số đi lên từ trái sang phải, nên hàm số đã cho đồng biến trên khoảng này
\item (Sai) Hàm số $f=|f(x)|$ nghịch biến trên khoảng $(-2,0)$. (Vì): Đồ thị hàm số $y=|f(x)|$ được suy ra từ đồ thị hàm số $y=f(x)$ bằng cách giữ nguyên phần đồ thị phía trên trục hoành và lấy đối xứng phần đồ thị phía dưới trục hoành qua trục hoành. Trên khoảng $(-2,0)$, đồ thị hàm số $y=f(x)$ nằm phía trên trục hoành và đang đi xuống (nghịch biến). Do đó, trên khoảng $(-2,0)$, hàm số $y=|f(x)|$ vẫn nghịch biến. Tuy nhiên, đề bài cho rằng hàm số $f=|f(x)|$ nghịch biến trên khoảng $(-2,0)$ là sai vì trên khoảng này $f(x)$ đang nghịch biến và $f(x)\gt 0$ nên $|f(x)|$ cũng nghịch biến. Mệnh đề này sai do cách diễn đạt "nghịch biến và đồng biến" ở lời giải gốc không rõ ràng và có thể gây nhầm lẫn. Xét kỹ hơn, trên khoảng $(-2,0)$, $f(x)$ nghịch biến và $f(x)\gt 0$, nên $|f(x)| = f(x)$ cũng nghịch biến. Tuy nhiên, nếu xét trên một khoảng rộng hơn bao gồm cả phần đồ thị âm, thì $|f(x)|$ có thể có cả tính đồng biến và nghịch biến. Với khoảng $(-2,0)$ cụ thể, hàm số $|f(x)|$ nghịch biến. Mệnh đề này sai do có sự nhầm lẫn trong việc phân tích đồ thị $|f(x)|$ hoặc do cách diễn đạt không chính xác
\end{itemchoice}
}
\end{ex}

% ===== Câu 2 =====
% ID: T12C1B1-01-TLN
% Mức: TH
\begin{ex}
% ID: T12C1B1-01-TLN
% Mức: TH
Cho hàm số $y=4x^3+8x^2-3x+3$. Giá trị cực đại của hàm số bằng bao nhiêu?
\shortans{12}
\loigiai{
Tập xác định $\mathbb{R}.\y'=12x^2+16x-3.\y'=0 \Leftrightarrow x=- \dfrac{3}{2} \text{ hoặc } x=\dfrac{1}{6}.\
  Hàm số đạt cực đại tạix=- \dfrac{3}{2}, y_{\text{CĐ}}=12$.
}
\end{ex}

% ===== Câu 3 =====
% ID: T12C1B1-01-TN
% Mức: TH
\begin{ex}
% ID: T12C1B1-01-TN
% Mức: TH
Cho hàm số $y=\dfrac{8-5x}{-4x-9}$. Khẳng định nào sau đây là $\textbf{sai}$?
\choice
{Hàm số đồng biến trên khoảng $\left(-\infty;-\dfrac{37}{4}\right)$}
{Hàm số đồng biến trên khoảng $\left(-\infty;-\dfrac{9}{4}\right)$}
{Hàm số đồng biến trên khoảng $\left(-\dfrac{9}{4};+\infty\right)$}
{\True Hàm số nghịch biến trên khoảng $\left(-\infty;-\dfrac{37}{4}\right)$}
\loigiai{
Tập xác định $\mathscr{D}=\mathbb{R}\setminus\left\{-\dfrac{9}{4}\right\}$.

Viết lại hàm số: $y=\dfrac{-5x+8}{-4x-9}$.

Ta có $y'=\dfrac{-5\cdot(-9)-8\cdot(-4)}{\left(-4x-9\right)^{2}}=\dfrac{77}{\left(-4x-9\right)^{2}}\gt 0,\forall x \in \mathscr{D}$.

Do đó, hàm số đồng biến trên mỗi khoảng $\left(-\infty;-\dfrac{9}{4}\right)$ và $\left(-\dfrac{9}{4};+\infty\right)$.

Vì $\left(-\infty;-\dfrac{37}{4}\right) \subset \left(-\infty;-\dfrac{9}{4}\right)$ nên hàm số cũng đồng biến trên khoảng $\left(-\infty;-\dfrac{37}{4}\right)$.
}
\end{ex}

% ===== Câu 4 =====
% ID: T12C1B1-02-DS
% Mức: TH
\begin{ex}
% ID: T12C1B1-02-DS
% Mức: TH
Cho hàm số $y=f(x)$ có đồ thị như hình vẽ bên dưới.
\choiceTF
{Hàm số đã cho chỉ có 2 cực trị}
{Hàm số đã cho có giá trị cực tiểu $x=-2$}
{\True Hàm số đã cho có giá trị cực đại $y=-1$}
{\True Đồ thị hàm số $y=|f(x)|$ có điểm cực đại là $(0,2)$}
\loigiai{
\begin{itemchoice}
\item (Đúng) Hàm số đã cho chỉ có 2 cực trị. (Vì): ựa vào đồ thị hàm số ta nhận thấy hàm số có 3 điểm cực trị
\item (Sai) Hàm số đã cho có giá trị cực tiểu $x=-2$. (Vì): Hàm số đã cho có giá trị cực tiểu $y=-2$
\item (Đúng) Hàm số đã cho có giá trị cực đại $y=-1$.
\item (Đúng) Đồ thị hàm số $y=|f(x)|$ có điểm cực đại là $(0,2)$. (Vì): Hàm số $y=|f(x)|$ có đồ thị như hình vẽ bên dưới
\end{itemchoice}
}
\end{ex}

% ===== Câu 5 =====
% ID: T12C1B1-02-TLN
% Mức: TH
\begin{ex}
% ID: T12C1B1-02-TLN
% Mức: TH
Tìm điểm cực tiểu của hàm số $y=x^3-5x^2-8x+8$?
\shortans{4}
\loigiai{
Tập xác định $\mathbb{R}$. $y' = 3x^2 - 10x - 8$. $y' = 0 \Leftrightarrow x = -\frac{2}{3}$ hoặc $x = 4$. Hàm số đạt cực tiểu tại $x = 4$, $y_{\text{CT}} = -40$.
}
\end{ex}

% ===== Câu 6 =====
% ID: T12C1B1-02-TN
% Mức: TH
\begin{ex}
% ID: T12C1B1-02-TN
% Mức: TH
Cho hàm số có bảng biến thiên như sau
 
 Khẳng định nào sau đây là \textbf{đúng}?
\choice
{Hàm số nghịch biến trên khoảng $\left(-\infty;-\dfrac{4}{3}\right)$}
{Hàm số đồng biến trên khoảng $\left(-\dfrac{4}{3};\dfrac{1}{2}\right)$}
{Hàm số nghịch biến trên khoảng $\left(\dfrac{7}{2};+\infty\right)$}
{\True Hàm số đồng biến trên khoảng $\left(-\infty;-\dfrac{7}{3}\right)$}
\loigiai{
Dựa vào bảng biến thiên ta thấy, hàm số đồng biến trên mỗi khoảng $\left(-\infty;-\dfrac{4}{3}\right)$ và $\left(\dfrac{1}{2};+\infty\right)$.
Hàm số nghịch biến trên $\left(-\dfrac{4}{3};\dfrac{1}{2}\right)$.
}
\end{ex}

% ===== Câu 7 =====
% ID: T12C1B1-03-TLN
% Mức: TH
\begin{ex}
% ID: T12C1B1-03-TLN
% Mức: TH
Cho hàm số $y=-x^3-mx^2+(4m+9)x+3$ với $m$ là tham số. Có bao nhiêu giá trị nguyên của $m$ để hàm số nghịch biến trên $(-\infty;+\infty)$?
\shortans{7}
\loigiai{
Tập xác định $\mathbb{R}$. $y'=-3x^2-2mx+4m+9$.
 Hàm số nghịch biến trên $(-\infty;+\infty)$ khi $a<0$ và $\Delta \leq 0$.
 $\Leftrightarrow m^2+12m+27 \leq 0\Leftrightarrow -9 \le m \le -3$.
 Suy ra $m \in \{-9,-8,-7,-6,-5,-4,-3\}$.
}
\end{ex}

% ===== Câu 8 =====
% ID: T12C1B1-03-TN
% Mức: TH
\begin{ex}
% ID: T12C1B1-03-TN
% Mức: TH
Cho hàm số có đồ thị như hình bên. Khẳng định nào sau đây là \textbf{đúng}?
\choice
{Hàm số nghịch biến trên khoảng $(-\infty;1)$}
{Hàm số nghịch biến trên khoảng $(-\infty;0)$}
{Hàm số nghịch biến trên khoảng $\left(\dfrac{7}{3};+\infty\right)$}
{\True Hàm số đồng biến trên khoảng $(-\infty;1)$}
\loigiai{
Dựa vào đồ thị ta thấy, hàm số đồng biến trên mỗi khoảng $(-\infty;1)$ và $\left(\dfrac{7}{3};+\infty\right)$. Hàm số nghịch biến trên $\left(1;\dfrac{7}{3}\right)$.
}
\end{ex}

% ===== Câu 9 =====
% ID: T12C1B1-04-TLN
% Mức: TH
\begin{ex}
% ID: T12C1B1-04-TLN
% Mức: TH
Đồ thị hàm số $y=x^3-3x^2+2ax+b$ có điểm cực tiểu $A(2;-2)$. Khi đó $a+b$ bằng bao nhiêu?
\shortans{2}
\loigiai{
$y' = 3x^2 - 6x + 2a$ (Điểm cực tiểu) $A(2;-2)$ thuộc đồ thị hàm số nên $y(2) = -2$.
 $2^3 - 3(2^2) + 2a(2) + b = -2 \Rightarrow 8 - 12 + 4a + b = -2 \Rightarrow 4a + b = -6$.
 Tại điểm cực tiểu, $y'(2) = 0$.
 $3(2^2) - 6(2) + 2a = 0 \Rightarrow 12 - 12 + 2a = 0 \Rightarrow 2a = 0 \Rightarrow a = 0$.
 Thay $a=0$ vào $4a+b=-6$, ta được $4(0) + b = -6 \Rightarrow b = -6$.
 $a+b = 0 + (-6) = -6$.
 -6
}
\end{ex}

% ===== Câu 10 =====
% ID: T12C1B1-04-TN
% Mức: TH
\begin{ex}
% ID: T12C1B1-04-TN
% Mức: TH
Tìm tất cả các giá trị thực của tham số $m$ để hàm số $y=-\dfrac{1}{3}x^3-mx^2+(2m-3)x-m+2$ nghịch biến trên $\mathbb{R}$.
\choice
{$m\le -3$, $m\ge 1$}
{$-3\lt m<1$}
{\True $-3\le m\le 1$}
{$m\le 1$}
\loigiai{
Ta có $y'=-x^2-2mx+(2m-3)$.
 Hàm số nghịch biến trên $\mathbb{R}$ khi
 $y'\le 0\Leftrightarrow \begin{cases} a=-1\lt 0 \ \Delta' \le 0 \end{cases}\Leftrightarrow m^2+2m-3\le 0\Leftrightarrow -3\le m \le 1$.
}
\end{ex}

% ===== Câu 11 =====
% ID: T12C1B1-05-TLN
% Mức: TH
\begin{ex}
% ID: T12C1B1-05-TLN
% Mức: TH
Có bao nhiêu giá trị nguyên của tham số $m$ để hàm số $y=x^3-3(m+2)x^2+3(m^2+4m)x+1$ nghịch biến trên khoảng $(0; 1)$?
\shortans{4}
\loigiai{
Hàm số nghịch biến trên $(0; 1)$ khi và chỉ khi $y'\leq 0,\quad \forall x\in (0; 1) \Leftrightarrow 3x^2-6(m+2)x+3(m^2+4m)\leq 0,\quad \forall x\in (0; 1) \Leftrightarrow (x-m)(3x-3m-12)\leq 0,\quad \forall x\in (0; 1) \Leftrightarrow m\leq x\leq m+4,\quad \forall x\in (0; 1) \Leftrightarrow \begin{cases} m\leq 0 \ 1\leq m+4 \end{cases} \Leftrightarrow -3\leq m\leq 0$. Suy ra có 4 giá trị của $m$ thỏa mãn đề bài.
}
\end{ex}

% ===== Câu 12 =====
% ID: T12C1B1-05-TN
% Mức: TH
\begin{ex}
% ID: T12C1B1-05-TN
% Mức: TH
Cho hàm số $y=-x^3-mx^2+(4m+9)x+5$, với $m$ là tham số thực. Có bao nhiêu giá trị nguyên của $m$ để hàm số nghịch biến trên $(-\infty;+\infty)$?
\choice
{$5$}
{$6$}
{\True $7$}
{$4$}
\loigiai{
Hàm số $y=-x^3-mx^2+(4m+9)x+5$ là hàm bậc 3 có hệ số $a=-1\lt 0$ nên điều kiện cần và đủ để $y=-x^3-mx^2+(4m+9)x+5$ nghịch biến trên $(-\infty;+\infty)$ là
  \begin{align*}
  &~ y'=-3x^2-2mx+ $4\,\text{m}$ +9\le 0, $\forall$ x $\in\mathbb{R}\Leftrightarrow$ &~ $\Delta$ '=m^2+ $12\,\text{m}$ +27\le 0 
\Leftrightarrow &~ -9\le m\le -3 
\Rightarrow &~ m $\in\{-9,-8,-7,-6,-5,-4,-3\}$.
  \end{align*}
}
\end{ex}

% ===== Câu 13 =====
% ID: T12C1B1-06-TLN
% Mức: TH
\begin{ex}
% ID: T12C1B1-06-TLN
% Mức: TH
Tìm giá trị của tham số $m$ để ba điểm cực trị của đồ thị hàm số $y=x^4+(6m-4)x^2+1-m$ là ba đỉnh của một tam giác vuông? (Kết quả làm tròn đến chữ số hàng phần trăm).
\shortans{0,33}
\loigiai{
Ta có $y'=4x^3+4(3m-2)x.$ Giải $y'=0\Leftrightarrow x=0$ hoặc $x^2=2-3m.$ (Hàm số có ba cực trị khi và chỉ khi) $y'=0$ có ba nghiệm phân biệt $\Leftrightarrow 2-3m>0\Leftrightarrow m<\dfrac{2}{3}.$ (Khi đó, tọa độ ba điểm cực trị là) $A(0;10m),B\left(\sqrt{2-3m};-(2-3m)^2+1-m\right)$ và $C\left(-\sqrt{2-3m};-(2-3m)^2+1-m\right).$ (Ta có) $\overrightarrow{AB}=\left(\sqrt{2-3m};-(2-3m)^2\right)$ và $\overrightarrow{AC}=\left(-\sqrt{2-3m};-(2-3m)^2\right).$ (Nhận xét:) $\triangle ABC$ luôn cân tại $A.$ (Do đó) $A$, $B$, $C$ tạo thành tam giác vuông $\Leftrightarrow \overrightarrow{AB} \cdot \overrightarrow{AC}=0\Leftrightarrow -(\sqrt{2-3m})^2+(-(2-3m)^2)^2=0\Leftrightarrow -(2-3m)+(2-3m)^4=0.\Leftrightarrow (2-3m)^4-(2-3m)=0\Leftrightarrow (2-3m)((2-3m)^3-1)=0.$ (Vì) $m<\dfrac{2}{3}$ nên $2-3m>0.$ Do đó $(2-3m)^3-1=0\Leftrightarrow (2-3m)^3=1\Leftrightarrow 2-3m=1\Leftrightarrow m=\dfrac{1}{3}.$
}
\end{ex}

% ===== Câu 14 =====
% ID: T12C1B1-06-TN
% Mức: TH
\begin{ex}
% ID: T12C1B1-06-TN
% Mức: TH
Cho hàm số $y=f(x)$ xác định trên $\mathbb{R}$ và có bảng biến thiên như hình vẽ sau
 
  
 
 Hàm số $y=f(x)$ đạt cực đại tại
\choice
{\True $x=10$}
{$x=8$}
{$x=12$}
{$x=17$}
\loigiai{
Dựa vào bảng biến thiên ta thấy hàm số đạt cực đại tại $x=10$.
}
\end{ex}

% ===== Câu 15 =====
% ID: T12C1B1-07-TN
% Mức: TH
\begin{ex}
% ID: T12C1B1-07-TN
% Mức: TH
Cho hàm số có đồ thị như hình bên. 
  Khẳng định nào $\textbf{sai}$?
\choice
{Hàm số có một cực tiểu tại $x=\dfrac{3}{4}$ và đạt cực đại tại $x=0$}
{Giá trị cực đại $y=-1$ và giá trị cực tiểu $y=-\dfrac{113}{32}$}
{Hàm số có một cực tiểu tại $x=-\dfrac{3}{4}$}
{\True Hàm số có một cực đại tại $x=-\dfrac{3}{4}$}
\loigiai{
Dựa vào đồ thị ta thấy, hàm số đạt cực đại tại $x=0, y_{\text{CĐ}}=-1$.
Hàm số đạt cực tiểu tại $x=-\dfrac{3}{4}$ và $x=\dfrac{3}{4}$, $y_{\text{CT}}=-\dfrac{113}{32}$.
}
\end{ex}

% ===== Câu 16 =====
% ID: T12C1B1-08-TN
% Mức: TH
\begin{ex}
% ID: T12C1B1-08-TN
% Mức: TH
Cho hàm số $y = 8x^3 -6mx^2 + 9mx + 6m$ ($m$ là tham số thực). Giá trị của $m$ để hàm số đạt cực tiểu tại $x=-6$ thuộc khoảng nào sau đây?
\choice
{$\left(-\infty;-\dfrac{35}{3}\right)$}
{$\left(-\infty;-\dfrac{47}{3}\right)$}
{$\left(-\dfrac{23}{3}; +\infty\right)$}
{\True $\varnothing$}
\loigiai{
Hàm số xác định với mọi $x \in \mathbb{R}.\
  Ta cóy'=24x^2 -12mx + 9m$.
Hàm số đạt cực tiểu tại $x=-6 \Rightarrow y'(-6) = 0 \Leftrightarrow 81m+864 = 0 \Leftrightarrow m=-\dfrac{32}{3}.\
  Thử lại:\
  Vớim = -\dfrac{32}{3}$, ta có:y'(-6) = 0và$y''(-6) = -160 \lt  0.\\Rightarrow x=-6$ là điểm cực đại của hàm số.
Vậy không có giá trị nào của $m$ thoả bài toán.
}
\end{ex}
